\section{实际系统中的并行策略}

\subsection{ZeRO：在数据并行中消除冗余}

DDP 的最大问题是每张卡都完整复制了参数、梯度和优化器状态。微软的 \textbf{ZeRO}（Zero Redundancy Optimizer）通过分片来消除这种冗余，分为三个递进的阶段：

\begin{knowledgebox}{ZeRO 的三个阶段}
\begin{itemize}
    \item \textbf{ZeRO-1}（优化器状态分片）：每个 rank 只保存 $\frac{1}{p}$ 的 Adam 状态（$m$ 和 $v$）。训练步骤中，每个 rank 只更新自己负责的那部分参数，然后通过 all-gather 广播更新后的参数。\textbf{显存节省}：以 13B 模型为例，Adam 状态从 104 GB/卡降到 26 GB/卡（4 卡）。
    \item \textbf{ZeRO-2}（+ 梯度分片）：在 ZeRO-1 基础上，梯度也只保留 $\frac{1}{p}$。反向传播时用 reduce-scatter 代替 all-reduce——每个 rank 只保留自己负责参数的归约梯度。\textbf{额外节省}：梯度从 26 GB/卡降到 6.5 GB/卡。
    \item \textbf{ZeRO-3}（+ 参数分片）：连参数本身也分片。前向和反向需要完整参数时，通过 all-gather 临时收集，用完立即释放。\textbf{极致节省}：每卡只需 $\frac{1}{p}$ 的全部训练状态。
\end{itemize}
\end{knowledgebox}

\begin{table}[H]
\centering
\small
\begin{tabular}{p{2.2cm}p{3cm}p{3cm}p{3.8cm}}
\toprule
\textbf{阶段} & \textbf{分片内容} & \textbf{每卡显存} & \textbf{额外通信} \\
\midrule
DDP（基线） & 无分片 & $16P$ bytes & 1 次 all-reduce \\
ZeRO-1 & 优化器状态 & $4P + \frac{12P}{p}$ bytes & 1 次 all-gather（参数） \\
ZeRO-2 & + 梯度 & $2P + \frac{14P}{p}$ bytes & reduce-scatter + all-gather \\
ZeRO-3 & + 参数 & $\frac{16P}{p}$ bytes & 前向/反向各 1 次 all-gather \\
\bottomrule
\end{tabular}
\caption{ZeRO 各阶段的显存与通信权衡。$P$ 为参数量，$p$ 为 GPU 数。}
\end{table}

\begin{warningbox}{ZeRO-3 的隐含代价}
ZeRO-3 虽然把显存降到最低，但它要求在前向和反向的每一层都做 all-gather 来收集完整参数——这和张量并行的通信频率类似。如果跨节点使用 ZeRO-3，通信开销可能反而比 ZeRO-2 + 流水线并行更大。实践中，ZeRO-2 通常是性价比最高的选择。
\end{warningbox}

\subsection{实际训练框架的并行策略对比}

\begin{table}[H]
\centering
\small
\resizebox{\textwidth}{!}{%
\begin{tabular}{p{3.0cm}p{3.0cm}p{3.0cm}p{3.0cm}p{3.0cm}}
\toprule
\textbf{维度} & \textbf{Megatron-LM} & \textbf{DeepSpeed ZeRO} & \textbf{PyTorch FSDP} \\
\midrule
核心思路 & 3D 并行 & 数据并行 + 状态分片 & 类 ZeRO-3 的参数分片 \\
张量并行 & 原生支持 & 需配合 Megatron & 不直接支持 \\
流水线并行 & 原生 (1F1B) & ZeRO + PP & 不直接支持 \\
数据并行 & 原生 & ZeRO-1/2/3 & 原生（分片式） \\
典型规模 & 100B+ 参数 & 10B--1T 参数 & 1B--100B 参数 \\
代码侵入性 & 高（需改模型） & 中（包装优化器） & 中（包装模块） \\
适用场景 & 超大规模预训练 & 灵活部署 & PyTorch 生态内 \\
\bottomrule
\end{tabular}
}
\caption{主流分布式训练框架对比}
\end{table}

\begin{knowledgebox}{Megatron-LM 的 3D 并行}
Megatron-LM 是 NVIDIA 开发的大模型训练框架，它同时使用三种并行：
\begin{itemize}
    \item \textbf{张量并行}（TP）：在同一节点的 8 张 NVLink 连接的 GPU 间切分层宽。
    \item \textbf{流水线并行}（PP）：跨节点切分层深度，用 interleaved 1F1B 调度。
    \item \textbf{数据并行}（DP）：在 TP$\times$PP 组之间复制，用 all-reduce 同步梯度。
\end{itemize}
例如训练 GPT-3 175B 时，一种典型配置是 TP=8（节点内），PP=8（跨 8 个节点），DP=16（16 组副本），共需 $8 \times 8 \times 16 = 1024$ 张 GPU。
\end{knowledgebox}

